\section{Background and Motivation}\label{sec:motivation}

This section addresses three major themes:
\emph{the necessity of placing nodes over multiple banks}
(Section~\ref{sec:mot_banks}), \emph{the requirement of a pipelined read-write
	system capable of handling data hazards gracefully} (Section~\ref{sec:mot_rw}),
and \emph{desirable properties of a placement function}
(Section~\ref{sec:mot_conflict}).

\subsection{A Review of Replication, Multi-Port Banks and Placement}\label{sec:mot_banks}

The requirement of reading many nodes in parallel can be met in three ways:
replicating the data structure over several banks, building banks with many
read and write ports, and placing the nodes across many single-port banks with
a dedicated placement function.

\paragraph{Replication} Replicating the structure is naturally conflict-free,
offering maximum parallelism at the cost of storage. A conflict-free traversal
that reads $d$ nodes at once needs $d$ copies, so a tree of $S$ bytes occupies
$S \cdot d$ bytes of on-chip memory, and every update must write all $d$
copies, so the write traffic grows with $d$. As a result, dense trees quickly
become infeasible.

\paragraph{Multi-port RAMs} Multi-port RAMs with 4, 8 or 16 read and write
ports could, in principle, sustain the bandwidth a traversal requires, provided
the number of ports $N$ is at least the number of nodes a traversal reads,
$N \ge d$. In a custom array, every port adds its own wordline and bitlines to
every cell, so the cell area rises roughly with the square of the number of
ports, and the longer, more heavily loaded lines raise access delay and
switching power with it~\cite{rixner2000}. An FPGA offers no such cells: its
block RAMs have two ports~\cite{ref2,ref30}, and more must be composed from
them, by multipumping at a multiple of the system clock, or by replication
steered by a live-value table, which builds $W$ write ports and $R$ read ports
from $W \cdot R$ replicas of the storage~\cite{laforest2010,laforest2014}, or
by XOR codes, which need $W(W - 1 + R)$ replicas~\cite{laforest2012}; later
refinements reduce the steering logic without changing the
replication~\cite{abdelhadi2014}. A tree traversal sits at the expensive end of
every one of these curves, because its writes are as wide as its reads.

\paragraph{Placement} Placing nodes into banks takes a different approach
altogether. A placement function $\operatorname{bank}(x)$ assigns each node $x$
to one of $B$ single-port banks, so that the tree is stored exactly once and
the nodes of every traversal land in distinct
banks~\cite{das2000mappings,zhou2017new}. For
a traversal to be conflict-free, $B$ must be at least the number of nodes it
reads from the banks, $B \ge d$, so the number of banks still grows with the
depth of the tree; but the storage stays at one copy, and each bank keeps a
single port, however deep the tree becomes. This is the approach we take.
Existing placements, however, keep either conflict freedom or balance, not both. Address
hashing~\cite{rau1991pseudorandomly,vandenbraak2016configurable} balances
capacity across banks, but it distributes addresses rather than separating the
ones a traversal requests together, so traversals still collide;
colouring~\cite{das2000mappings} separates co-accessed nodes by construction,
but gives up balance. Dynamic load-store queues and out-of-order
arbiters~\cite{josipovic2017outoforder} resolve conflicts at run time, at an
area and frequency cost~\cite{szafarczyk2023compiler} that is warranted for
unpredictable access, but not for a traversal whose access set is a function of
its index.

% rw_throughput_fig.tex: cycles per operation of a pipeline that stalls on read-after-write hazards, over the
% whole ITCH trace; greyscale. Self-contained (data inline).
% Needs: \usepackage{pgfplots} \pgfplotsset{compat=1.17} 
\usepgfplotslibrary{fillbetween}

\pgfplotstableread{
	x med lo hi
	0.0050 1.3080 1.2436 1.3524
	0.0150 1.3160 1.2544 1.3720
	0.0250 1.3220 1.2592 1.4240
	0.0350 1.3120 1.2520 1.3764
	0.0450 1.2680 1.2108 1.3252
	0.0550 1.3000 1.2144 1.3320
	0.0650 1.3140 1.2396 1.3720
	0.0750 1.3320 1.2832 1.3680
	0.0850 1.3040 1.2680 1.3632
	0.0950 1.3020 1.2556 1.3484
	0.1050 1.3080 1.2236 1.3560
	0.1150 1.3100 1.2500 1.3764
	0.1250 1.3060 1.2600 1.3680
	0.1350 1.2680 1.1996 1.3404
	0.1450 1.2860 1.2316 1.3280
	0.1550 1.3080 1.2356 1.3564
	0.1650 1.3380 1.2676 1.4040
	0.1750 1.3040 1.2200 1.3648
	0.1850 1.3140 1.2316 1.3732
	0.1950 1.2760 1.2112 1.3208
	0.2050 1.2980 1.2360 1.3640
	0.2150 1.3080 1.2236 1.3604
	0.2250 1.3280 1.2800 1.3720
	0.2350 1.3220 1.2836 1.3768
	0.2450 1.3180 1.2760 1.3844
	0.2550 1.3100 1.2636 1.3688
	0.2650 1.3240 1.2636 1.3760
	0.2750 1.3240 1.2796 1.3720
	0.2850 1.3040 1.2712 1.3480
	0.2950 1.3240 1.2600 1.3720
	0.3050 1.3100 1.2756 1.3884
	0.3150 1.3060 1.2664 1.3604
	0.3250 1.3140 1.2468 1.3772
	0.3350 1.3120 1.2540 1.3760
	0.3450 1.3160 1.2632 1.3648
	0.3550 1.3200 1.2312 1.3804
	0.3650 1.3160 1.2632 1.3688
	0.3750 1.3240 1.2712 1.3684
	0.3850 1.3480 1.3000 1.3804
	0.3950 1.3180 1.2756 1.3688
	0.4050 1.3060 1.2220 1.3456
	0.4150 1.3120 1.2556 1.3640
	0.4250 1.3140 1.2792 1.3564
	0.4350 1.3340 1.2584 1.3968
	0.4450 1.3200 1.2556 1.3604
	0.4550 1.3160 1.2076 1.3760
	0.4650 1.3180 1.2596 1.3684
	0.4750 1.3160 1.2472 1.3680
	0.4850 1.3300 1.2756 1.3760
	0.4950 1.3060 1.2600 1.3620
	0.5050 1.3160 1.2716 1.3640
	0.5150 1.3260 1.2520 1.3772
	0.5250 1.3240 1.2320 1.3648
	0.5350 1.3080 1.2472 1.3564
	0.5450 1.3180 1.2480 1.3688
	0.5550 1.3000 1.2588 1.3600
	0.5650 1.3200 1.2560 1.3524
	0.5750 1.3220 1.2716 1.3804
	0.5850 1.3280 1.2916 1.3764
	0.5950 1.3280 1.2676 1.3804
	0.6050 1.3360 1.2676 1.3768
	0.6150 1.3280 1.2756 1.3920
	0.6250 1.3280 1.2840 1.3660
	0.6350 1.2960 1.2428 1.3524
	0.6450 1.3320 1.2596 1.3768
	0.6550 1.3380 1.2612 1.4012
	0.6650 1.3180 1.2556 1.3680
	0.6750 1.3160 1.2672 1.3768
	0.6850 1.3400 1.2840 1.3884
	0.6950 1.3080 1.2596 1.3528
	0.7050 1.3200 1.2712 1.3684
	0.7150 1.3140 1.2712 1.3688
	0.7250 1.3000 1.2300 1.3644
	0.7350 1.2920 1.2192 1.3284
	0.7450 1.2980 1.2668 1.3488
	0.7550 1.3020 1.2672 1.3600
	0.7650 1.3000 1.2520 1.3480
	0.7750 1.3140 1.2620 1.3600
	0.7850 1.3080 1.2420 1.3684
	0.7950 1.2940 1.2036 1.3800
	0.8050 1.3260 1.2556 1.3724
	0.8150 1.3320 1.2664 1.3724
	0.8250 1.3280 1.2552 1.3844
	0.8350 1.3360 1.2672 1.4124
	0.8450 1.3440 1.2624 1.4004
	0.8550 1.3020 1.2356 1.3844
	0.8650 1.3160 1.2520 1.3764
	0.8750 1.2940 1.2480 1.3604
	0.8850 1.2920 1.2556 1.3488
	0.8950 1.2960 1.2200 1.3528
	0.9050 1.3200 1.2464 1.3844
	0.9150 1.3120 1.2272 1.3680
	0.9250 1.3020 1.2236 1.3340
	0.9350 1.2680 1.2080 1.3200
	0.9450 1.3120 1.2276 1.3448
	0.9550 1.2860 1.2072 1.3368
	0.9650 1.2880 1.2228 1.3332
	0.9750 1.2860 1.2152 1.3692
	0.9850 1.3080 1.2264 1.3588
	0.9950 1.3100 1.2396 1.3644
	1.0050 1.3040 1.2516 1.3604
	1.0150 1.3020 1.2308 1.3420
	1.0250 1.2920 1.2620 1.3488
	1.0350 1.3160 1.2596 1.3808
	1.0450 1.3100 1.2356 1.3680
	1.0550 1.3060 1.2588 1.3612
	1.0650 1.3020 1.2600 1.3484
	1.0750 1.3180 1.2224 1.3804
	1.0850 1.2940 1.1996 1.3768
	1.0950 1.3420 1.2720 1.3932
	1.1050 1.3060 1.2464 1.3644
	1.1150 1.3180 1.2476 1.3920
	1.1250 1.3160 1.2236 1.3804
	1.1350 1.3060 1.2480 1.3564
	1.1450 1.3160 1.2460 1.3684
	1.1550 1.3400 1.2796 1.3808
	1.1650 1.3160 1.2872 1.3964
	1.1750 1.3200 1.2756 1.3748
	1.1850 1.3000 1.2436 1.3692
	1.1950 1.3180 1.2548 1.3848
	1.2050 1.3240 1.2860 1.3808
	1.2150 1.3180 1.2636 1.3724
	1.2250 1.3120 1.2436 1.3844
	1.2350 1.3100 1.2356 1.3600
	1.2450 1.3260 1.2756 1.3844
	1.2550 1.3080 1.2160 1.3484
	1.2650 1.3020 1.2160 1.3564
	1.2750 1.3220 1.2512 1.3760
	1.2850 1.3260 1.2700 1.3848
	1.2950 1.3120 1.2508 1.3612
	1.3050 1.3160 1.2712 1.3644
	1.3150 1.3080 1.2388 1.3844
	1.3250 1.3120 1.2620 1.3720
	1.3350 1.3060 1.2596 1.3640
	1.3450 1.3280 1.2680 1.3604
	1.3550 1.3060 1.2348 1.3764
	1.3650 1.3260 1.2548 1.3760
	1.3750 1.2920 1.2240 1.3440
	1.3850 1.2960 1.2428 1.3600
	1.3950 1.3040 1.2400 1.3564
	1.4050 1.2960 1.2184 1.3480
	1.4150 1.3080 1.2312 1.3560
	1.4250 1.2540 1.1992 1.3160
	1.4350 1.2700 1.2228 1.3532
	1.4450 1.2980 1.2400 1.3532
	1.4550 1.2900 1.2152 1.3632
	1.4650 1.2840 1.1840 1.3408
	1.4750 1.3000 1.2480 1.3812
	1.4850 1.3220 1.2156 1.3924
	1.4950 1.3120 1.2280 1.3840
	1.5050 1.2880 1.2196 1.3444
	1.5150 1.3040 1.2432 1.3684
	1.5250 1.3140 1.2508 1.3780
	1.5350 1.3080 1.2512 1.3964
	1.5450 1.2860 1.2108 1.3576
	1.5550 1.2760 1.2160 1.3640
	1.5650 1.3260 1.2560 1.3800
	1.5750 1.3140 1.2304 1.3920
	1.5850 1.3240 1.2396 1.3880
	1.5950 1.3280 1.2632 1.3924
	1.6050 1.3120 1.2476 1.3808
	1.6150 1.3400 1.2316 1.3964
	1.6250 1.3340 1.2192 1.3768
	1.6350 1.3100 1.2196 1.3560
	1.6450 1.3200 1.2624 1.3724
	1.6550 1.3480 1.2476 1.3812
	1.6650 1.3080 1.2356 1.3444
	1.6750 1.3020 1.2396 1.3564
	1.6850 1.3360 1.2588 1.3680
	1.6950 1.3060 1.2392 1.3644
	1.7050 1.3460 1.2640 1.3728
	1.7150 1.3080 1.2676 1.3772
	1.7250 1.3040 1.2552 1.3648
	1.7350 1.3380 1.2872 1.4084
	1.7450 1.3140 1.2600 1.3564
	1.7550 1.3360 1.2828 1.3848
	1.7650 1.3360 1.2920 1.3764
	1.7750 1.3100 1.2356 1.3692
	1.7850 1.3260 1.2512 1.3728
	1.7950 1.3220 1.2552 1.3728
	1.8050 1.3080 1.2596 1.3804
	1.8150 1.3160 1.2392 1.3964
	1.8250 1.3040 1.2476 1.3528
	1.8350 1.3000 1.2588 1.3480
	1.8450 1.3100 1.2600 1.3768
	1.8550 1.3100 1.2708 1.3520
	1.8650 1.3200 1.2476 1.3648
	1.8750 1.3100 1.2620 1.3644
	1.8850 1.3080 1.2512 1.3404
	1.8950 1.3120 1.2672 1.3484
	1.9050 1.3200 1.2676 1.3644
	1.9150 1.3020 1.2640 1.3524
	1.9250 1.3080 1.2556 1.3540
	1.9350 1.3360 1.2752 1.3764
	1.9450 1.3000 1.2556 1.3640
	1.9550 1.3200 1.2516 1.3728
	1.9650 1.3120 1.2396 1.3684
	1.9750 1.3040 1.2440 1.3564
	1.9850 1.3080 1.2584 1.3532
	1.9950 1.3040 1.2480 1.3680
	2.0050 1.3160 1.2556 1.3564
	2.0150 1.3240 1.2672 1.3724
	2.0250 1.3100 1.2276 1.3684
	2.0350 1.2920 1.2396 1.3528
	2.0450 1.2760 1.2396 1.3408
	2.0550 1.3020 1.2592 1.3548
	2.0650 1.3080 1.2280 1.3644
	2.0750 1.2880 1.2312 1.3520
	2.0850 1.3120 1.2352 1.3532
	2.0950 1.3040 1.2236 1.3488
	2.1050 1.3240 1.2548 1.3700
	2.1150 1.3160 1.2796 1.3764
	2.1250 1.2960 1.2520 1.3564
	2.1350 1.2760 1.1796 1.3524
	2.1450 1.3080 1.2420 1.3768
	2.1550 1.2940 1.2272 1.3684
	2.1650 1.3240 1.2420 1.3852
	2.1750 1.3040 1.2640 1.3728
	2.1850 1.3040 1.2520 1.3368
	2.1950 1.3060 1.2496 1.3644
	2.2050 1.2940 1.2036 1.3400
	2.2150 1.2840 1.2312 1.3488
	2.2250 1.2960 1.2472 1.3444
	2.2350 1.3220 1.2628 1.3644
	2.2450 1.3240 1.2440 1.3844
	2.2550 1.3360 1.2588 1.3928
	2.2650 1.3080 1.2472 1.3640
	2.2750 1.3100 1.2472 1.3564
	2.2850 1.2780 1.2072 1.3528
	2.2950 1.2360 1.1676 1.3520
	2.3050 1.3000 1.2148 1.3564
	2.3150 1.2920 1.2236 1.3564
	2.3250 1.2880 1.2476 1.3640
	2.3350 1.3000 1.2320 1.3452
	2.3450 1.3140 1.2624 1.3608
	2.3550 1.2980 1.2116 1.3644
	2.3650 1.2780 1.2276 1.3424
	2.3750 1.2860 1.2188 1.3600
	2.3850 1.2980 1.2032 1.3800
	2.3950 1.2820 1.2276 1.3920
	2.4050 1.2840 1.2156 1.3492
	2.4150 1.3000 1.2420 1.3564
	2.4250 1.3000 1.2352 1.3732
	2.4350 1.3120 1.2392 1.3728
	2.4450 1.3000 1.2276 1.3520
	2.4550 1.3020 1.2392 1.3692
	2.4650 1.3100 1.2476 1.3600
	2.4750 1.3140 1.2612 1.3568
	2.4850 1.2980 1.2276 1.3520
	2.4950 1.3020 1.2160 1.3604
	2.5050 1.3360 1.2680 1.3800
	2.5150 1.3360 1.2672 1.3844
	2.5250 1.2980 1.2276 1.3780
	2.5350 1.2860 1.1876 1.3564
	2.5450 1.2760 1.2304 1.3604
	2.5550 1.3100 1.2272 1.3404
	2.5650 1.3280 1.2840 1.3808
	2.5750 1.2800 1.1988 1.3492
	2.5850 1.2980 1.2116 1.3688
	2.5950 1.2860 1.2232 1.3452
	2.6050 1.3000 1.2344 1.3640
	2.6150 1.2880 1.2276 1.3568
	2.6250 1.3020 1.2392 1.3600
	2.6350 1.3100 1.2636 1.3884
	2.6450 1.2720 1.2152 1.3216
	2.6550 1.2920 1.2036 1.3564
	2.6650 1.3120 1.2348 1.3680
	2.6750 1.2900 1.2312 1.3604
	2.6850 1.2900 1.2476 1.3408
	2.6950 1.3020 1.2280 1.3652
	2.7050 1.2640 1.1712 1.3480
	2.7150 1.3080 1.2556 1.3520
	2.7250 1.2680 1.2152 1.3324
	2.7350 1.3160 1.2600 1.3760
	2.7450 1.3020 1.2196 1.3692
	2.7550 1.3060 1.2436 1.3644
	2.7650 1.3100 1.2308 1.3612
	2.7750 1.2860 1.2036 1.3560
	2.7850 1.2740 1.2192 1.3524
	2.7950 1.2800 1.2276 1.3484
	2.8050 1.2660 1.2152 1.3440
	2.8150 1.2980 1.2116 1.3732
	2.8250 1.2960 1.2088 1.3536
	2.8350 1.2960 1.2296 1.3484
	2.8450 1.3040 1.2292 1.3568
	2.8550 1.2920 1.1952 1.3560
	2.8650 1.2880 1.2068 1.3960
	2.8750 1.2740 1.2200 1.3412
	2.8850 1.2860 1.2320 1.3568
	2.8950 1.2920 1.2356 1.3644
	2.9050 1.2920 1.1936 1.3724
	2.9150 1.2840 1.2156 1.3796
	2.9250 1.3040 1.2380 1.3564
	2.9350 1.2740 1.2160 1.3412
	2.9450 1.2800 1.2180 1.3520
	2.9550 1.2960 1.2512 1.3404
	2.9650 1.2740 1.2316 1.3724
	2.9750 1.2720 1.2240 1.3284
	2.9850 1.3020 1.2476 1.3684
	2.9950 1.2840 1.2388 1.3244
	3.0050 1.2960 1.2472 1.3484
	3.0150 1.2960 1.2316 1.3536
	3.0250 1.3080 1.2360 1.3724
	3.0350 1.2720 1.2028 1.3336
	3.0450 1.2720 1.1944 1.3244
	3.0550 1.3100 1.2260 1.3840
	3.0650 1.2560 1.1596 1.3268
	3.0750 1.2980 1.2096 1.3644
	3.0850 1.2920 1.2352 1.3736
	3.0950 1.3080 1.2468 1.3724
	3.1050 1.2960 1.2496 1.3496
	3.1150 1.3020 1.2516 1.3648
	3.1250 1.2860 1.2192 1.3444
	3.1350 1.2880 1.1996 1.3808
	3.1450 1.2900 1.2356 1.3604
	3.1550 1.2860 1.2100 1.3532
	3.1650 1.3020 1.2232 1.3724
	3.1750 1.3240 1.2796 1.3644
	3.1850 1.3100 1.2612 1.3688
	3.1950 1.3380 1.2744 1.3924
	3.2050 1.3060 1.2596 1.3720
	3.2150 1.3320 1.2860 1.4048
	3.2250 1.3320 1.2872 1.3928
	3.2350 1.3580 1.2988 1.4000
	3.2450 1.3320 1.2784 1.3808
	3.2550 1.3280 1.2596 1.3964
	3.2650 1.3260 1.2360 1.3880
	3.2750 1.3200 1.2664 1.3668
	3.2850 1.3000 1.2352 1.3764
	3.2950 1.2960 1.2400 1.3440
	3.3050 1.3420 1.2716 1.3924
	3.3150 1.3600 1.2840 1.4044
	3.3250 1.3340 1.2508 1.3976
	3.3350 1.3120 1.2600 1.3560
	3.3450 1.3400 1.2868 1.3768
	3.3550 1.3340 1.2600 1.3760
	3.3650 1.3420 1.2760 1.3764
	3.3750 1.3260 1.2408 1.4080
	3.3850 1.3020 1.1588 1.3760
	3.3950 1.2680 1.1552 1.3884
	3.4050 1.3040 1.1716 1.4048
	3.4150 1.3300 1.2400 1.4360
	3.4250 1.3920 1.1428 1.4844
	3.4350 1.4780 1.3828 1.4964
	3.4450 1.3060 1.2680 1.3724
	3.4550 1.3460 1.2792 1.3920
	3.4650 1.3580 1.3152 1.4084
	3.4750 1.3500 1.2836 1.4004
	3.4850 1.3480 1.3064 1.4204
	3.4950 1.3440 1.3076 1.3924
	3.5050 1.3760 1.3144 1.4088
	3.5150 1.3380 1.2880 1.4000
	3.5250 1.3620 1.3068 1.4124
	3.5350 1.3520 1.3000 1.4080
	3.5450 1.3560 1.2800 1.4004
	3.5550 1.3520 1.2960 1.3804
	3.5650 1.3540 1.2980 1.3964
	3.5750 1.3560 1.3068 1.4040
	3.5850 1.3660 1.2940 1.4164
	3.5950 1.3560 1.3064 1.4004
	3.6050 1.3460 1.2836 1.3888
	3.6150 1.3600 1.3228 1.4004
	3.6250 1.3500 1.2840 1.4200
	3.6350 1.3400 1.2676 1.3844
	3.6450 1.3600 1.3028 1.4088
	3.6550 1.3580 1.3116 1.3964
	3.6650 1.3480 1.2908 1.3960
	3.6750 1.3620 1.2868 1.4124
	3.6850 1.3620 1.3000 1.4080
	3.6950 1.3480 1.2872 1.4136
	3.7050 1.3360 1.2824 1.3844
	3.7150 1.3640 1.2960 1.4164
	3.7250 1.3780 1.3280 1.4212
	3.7350 1.3640 1.3112 1.3960
	3.7450 1.3600 1.3040 1.4044
	3.7550 1.3560 1.2880 1.4004
	3.7650 1.3680 1.2996 1.4124
	3.7750 1.3600 1.2872 1.3964
	3.7850 1.3440 1.2752 1.3960
	3.7950 1.3840 1.3116 1.4280
	3.8050 1.3800 1.3148 1.4124
	3.8150 1.3720 1.3192 1.3960
	3.8250 1.3500 1.2756 1.4040
	3.8350 1.3340 1.2940 1.4000
	3.8450 1.3480 1.3040 1.4080
	3.8550 1.3680 1.3116 1.4128
	3.8650 1.3640 1.3072 1.4044
	3.8750 1.3600 1.3076 1.4124
	3.8850 1.3440 1.2996 1.3928
	3.8950 1.3540 1.3104 1.4120
	3.9050 1.3540 1.3076 1.4004
	3.9150 1.3460 1.3200 1.3888
	3.9250 1.3340 1.2920 1.4000
	3.9350 1.3640 1.2912 1.4044
	3.9450 1.3480 1.3076 1.3840
	3.9550 1.3520 1.2988 1.3932
	3.9650 1.3640 1.3076 1.4164
	3.9750 1.3620 1.3080 1.4084
	3.9850 1.3500 1.3036 1.3880
	3.9950 1.3420 1.2912 1.3924
	4.0050 1.3660 1.3232 1.3976
	4.0150 1.3680 1.2992 1.4168
	4.0250 1.3600 1.2868 1.4044
	4.0350 1.3540 1.2952 1.4004
	4.0450 1.3720 1.3076 1.4160
	4.0550 1.3620 1.2876 1.4016
	4.0650 1.3700 1.2912 1.4092
	4.0750 1.3700 1.3248 1.4168
	4.0850 1.3580 1.2636 1.4040
	4.0950 1.3440 1.2956 1.4120
	4.1050 1.3280 1.2836 1.3804
	4.1150 1.3420 1.2904 1.3856
	4.1250 1.3540 1.3032 1.4084
	4.1350 1.3680 1.3196 1.4020
	4.1450 1.3620 1.3036 1.4280
	4.1550 1.3760 1.3308 1.4204
	4.1650 1.3620 1.3072 1.4008
	4.1750 1.3600 1.3152 1.4044
	4.1850 1.3580 1.2840 1.3924
	4.1950 1.3700 1.3160 1.4288
	4.2050 1.3700 1.3240 1.4164
	4.2150 1.3720 1.3316 1.4120
	4.2250 1.4040 1.3480 1.4368
	4.2350 1.3860 1.3512 1.4208
	4.2450 1.3720 1.3236 1.4168
	4.2550 1.3840 1.3152 1.4160
	4.2650 1.3700 1.3108 1.4124
	4.2750 1.3600 1.3188 1.4084
	4.2850 1.3640 1.3144 1.4044
	4.2950 1.3740 1.3196 1.4120
	4.3050 1.3480 1.2960 1.3924
	4.3150 1.3700 1.3196 1.4124
	4.3250 1.3800 1.3436 1.4204
	4.3350 1.3640 1.3144 1.4216
	4.3450 1.3700 1.3116 1.4088
	4.3550 1.3660 1.3192 1.4088
	4.3650 1.3860 1.3276 1.4244
	4.3750 1.3620 1.3116 1.4208
	4.3850 1.3700 1.3320 1.4168
	4.3950 1.3620 1.2832 1.4120
	4.4050 1.3520 1.3160 1.4168
	4.4150 1.3640 1.3232 1.3968
	4.4250 1.3780 1.3236 1.4172
	4.4350 1.3640 1.3156 1.3972
	4.4450 1.3720 1.3276 1.4244
	4.4550 1.3620 1.3156 1.4164
	4.4650 1.3900 1.3320 1.4084
	4.4750 1.3880 1.3472 1.4248
	4.4850 1.3760 1.3276 1.4124
	4.4950 1.3820 1.3112 1.4284
	4.5050 1.3760 1.3080 1.4244
	4.5150 1.3380 1.0532 1.4248
	4.5250 1.3760 1.3276 1.4124
	4.5350 1.3760 1.3216 1.4120
	4.5450 1.3760 1.3264 1.4168
	4.5550 1.3820 1.3200 1.4324
	4.5650 1.3800 1.3276 1.4240
	4.5750 1.3760 1.3356 1.4168
	4.5850 1.3780 1.3356 1.4244
	4.5950 1.3780 1.3320 1.4084
	4.6050 1.3720 1.3256 1.4124
	4.6150 1.3640 1.3116 1.3984
	4.6250 1.3540 1.3148 1.3932
	4.6350 1.3840 1.3104 1.4212
	4.6450 1.3660 1.3168 1.4204
	4.6550 1.3580 1.3280 1.3964
	4.6650 1.3600 1.3232 1.4052
	4.6750 1.3820 1.3396 1.4164
	4.6850 1.3680 1.3236 1.4084
	4.6950 1.3720 1.3120 1.4120
	4.7050 1.3640 1.3080 1.4124
	4.7150 1.3620 1.3232 1.4124
	4.7250 1.3640 1.3028 1.4084
	4.7350 1.3660 1.3120 1.4120
	4.7450 1.3580 1.3080 1.3936
	4.7550 1.3560 1.3040 1.4124
	4.7650 1.3580 1.3076 1.4040
	4.7750 1.3380 1.3000 1.3884
	4.7850 1.3620 1.3024 1.4048
	4.7950 1.3720 1.3304 1.4044
	4.8050 1.3540 1.2908 1.4040
	4.8150 1.3620 1.3076 1.3976
	4.8250 1.3580 1.3016 1.4084
	4.8350 1.3620 1.3356 1.4004
	4.8450 1.3800 1.3356 1.4088
	4.8550 1.3800 1.3120 1.4244
	4.8650 1.3640 1.3156 1.4008
	4.8750 1.3740 1.2920 1.4244
	4.8850 1.3620 1.3040 1.3924
	4.8950 1.3600 1.3080 1.4000
	4.9050 1.3560 1.3160 1.3924
	4.9150 1.3500 1.3156 1.4044
	4.9250 1.3720 1.3156 1.4204
	4.9350 1.3660 1.3152 1.4048
	4.9450 1.3840 1.3240 1.4044
	4.9550 1.3720 1.3196 1.4172
	4.9650 1.3680 1.3160 1.4088
	4.9750 1.3660 1.3116 1.4164
	4.9850 1.3600 1.3024 1.3960
	4.9950 1.3480 1.2760 1.4084
	5.0050 1.3680 1.2932 1.4208
	5.0150 1.3740 1.2996 1.4120
	5.0250 1.3540 1.2976 1.4040
	5.0350 1.3460 1.2828 1.4048
	5.0450 1.3740 1.3116 1.4000
	5.0550 1.3500 1.2904 1.4016
	5.0650 1.3520 1.3160 1.4160
	5.0750 1.3660 1.3340 1.4088
	5.0850 1.3840 1.3432 1.4280
	5.0950 1.3520 1.2956 1.4056
	5.1050 1.3660 1.3360 1.3964
	5.1150 1.3600 1.2868 1.4052
	5.1250 1.3600 1.2908 1.4004
	5.1350 1.3320 1.2704 1.3824
	5.1450 1.3560 1.3120 1.4048
	5.1550 1.3740 1.3160 1.4164
	5.1650 1.3480 1.2840 1.3848
	5.1750 1.3480 1.2832 1.4084
	5.1850 1.3380 1.2748 1.3964
	5.1950 1.3600 1.2992 1.4128
	5.2050 1.3620 1.3080 1.4048
	5.2150 1.3380 1.2904 1.4048
	5.2250 1.3500 1.2952 1.4044
	5.2350 1.3400 1.2880 1.4048
	5.2450 1.3540 1.2996 1.4088
	5.2550 1.3680 1.3188 1.4088
	5.2650 1.3540 1.2996 1.4044
	5.2750 1.3600 1.2996 1.3948
	5.2850 1.3580 1.3036 1.4000
	5.2950 1.3540 1.2832 1.4040
	5.3050 1.3600 1.3036 1.4128
	5.3150 1.3740 1.3120 1.4124
	5.3250 1.3580 1.2996 1.4160
	5.3350 1.3700 1.3000 1.4204
	5.3450 1.3860 1.2996 1.4280
	5.3550 1.3600 1.3040 1.4120
	5.3650 1.3760 1.3000 1.4160
	5.3750 1.3460 1.3068 1.4044
	5.3850 1.3640 1.3212 1.4056
	5.3950 1.3760 1.2996 1.4084
	5.4050 1.3680 1.3032 1.4208
	5.4150 1.3180 1.2748 1.3880
	5.4250 1.3580 1.2720 1.4044
	5.4350 1.3500 1.3232 1.4244
	5.4450 1.3520 1.2824 1.4044
	5.4550 1.3580 1.2752 1.4240
	5.4650 1.3460 1.2956 1.4052
	5.4750 1.3600 1.3060 1.4164
	5.4850 1.3660 1.3036 1.4088
	5.4950 1.3420 1.2880 1.4128
	5.5050 1.3640 1.3036 1.4004
	5.5150 1.3600 1.3092 1.3964
	5.5250 1.3380 1.2912 1.3844
	5.5350 1.3420 1.2952 1.3840
	5.5450 1.3580 1.2920 1.4124
	5.5550 1.3320 1.2916 1.3884
	5.5650 1.3720 1.3152 1.4124
	5.5750 1.3420 1.2952 1.4040
	5.5850 1.3280 1.2716 1.3884
	5.5950 1.3660 1.2716 1.4124
	5.6050 1.3480 1.2756 1.3928
	5.6150 1.3240 1.2632 1.3932
	5.6250 1.3520 1.2872 1.4128
	5.6350 1.3600 1.2996 1.4240
	5.6450 1.3560 1.3032 1.4044
	5.6550 1.3480 1.2920 1.3964
	5.6650 1.3580 1.2836 1.4040
	5.6750 1.3380 1.2960 1.3892
	5.6850 1.3260 1.2556 1.4008
	5.6950 1.3720 1.3120 1.4004
	5.7050 1.3460 1.2920 1.3960
	5.7150 1.3480 1.2880 1.3936
	5.7250 1.3360 1.2760 1.3960
	5.7350 1.3380 1.2908 1.4044
	5.7450 1.3540 1.3000 1.3884
	5.7550 1.3540 1.2840 1.4084
	5.7650 1.3680 1.3192 1.4040
	5.7750 1.3480 1.2760 1.4092
	5.7850 1.3260 1.2584 1.3780
	5.7950 1.3520 1.3080 1.4120
	5.8050 1.3420 1.2800 1.4040
	5.8150 1.3360 1.2796 1.3924
	5.8250 1.3360 1.2632 1.3964
	5.8350 1.3280 1.2632 1.3668
	5.8450 1.3340 1.2756 1.3968
	5.8550 1.3420 1.2708 1.4084
	5.8650 1.3440 1.2832 1.3800
	5.8750 1.3660 1.2992 1.4004
	5.8850 1.3440 1.2876 1.4004
	5.8950 1.3560 1.3156 1.4080
	5.9050 1.3360 1.3036 1.4160
	5.9150 1.3420 1.2904 1.3920
	5.9250 1.3380 1.2708 1.3968
	5.9350 1.3300 1.2640 1.4124
	5.9450 1.3260 1.2676 1.3764
	5.9550 1.3420 1.2908 1.3884
	5.9650 1.3420 1.2912 1.3980
	5.9750 1.3480 1.2912 1.4168
	5.9850 1.3520 1.3160 1.3924
	5.9950 1.3020 1.2584 1.3920
	6.0050 1.3500 1.3188 1.4084
	6.0150 1.3320 1.2688 1.3840
	6.0250 1.3360 1.2868 1.4160
	6.0350 1.3520 1.3072 1.4124
	6.0450 1.3500 1.3076 1.4084
	6.0550 1.3700 1.3112 1.4136
	6.0650 1.3500 1.3112 1.4084
	6.0750 1.3440 1.2664 1.3884
	6.0850 1.3440 1.2840 1.3848
	6.0950 1.3660 1.2916 1.3968
	6.1050 1.3420 1.2760 1.3848
	6.1150 1.3240 1.2516 1.3720
	6.1250 1.3320 1.2676 1.3808
	6.1350 1.3360 1.2800 1.3896
	6.1450 1.3520 1.2828 1.4016
	6.1550 1.3400 1.2912 1.4040
	6.1650 1.3440 1.2920 1.4084
	6.1750 1.3440 1.3040 1.3880
	6.1850 1.3280 1.2600 1.3844
	6.1950 1.3400 1.2596 1.3892
	6.2050 1.3460 1.2864 1.3964
	6.2150 1.3460 1.3000 1.3924
	6.2250 1.3400 1.2788 1.3960
	6.2350 1.3480 1.2756 1.3976
	6.2450 1.3320 1.2676 1.3720
	6.2550 1.3660 1.3032 1.4124
	6.2650 1.3520 1.2796 1.3968
	6.2750 1.3460 1.2556 1.3844
	6.2850 1.3160 1.2636 1.3764
	6.2950 1.3580 1.3188 1.4044
	6.3050 1.3500 1.2868 1.4128
	6.3150 1.3560 1.3072 1.4008
	6.3250 1.3500 1.3036 1.3920
	6.3350 1.3540 1.2952 1.4140
	6.3450 1.3640 1.2980 1.4240
	6.3550 1.3500 1.2876 1.4204
	6.3650 1.3600 1.3120 1.4164
	6.3750 1.3560 1.2952 1.4092
	6.3850 1.3680 1.2920 1.4164
	6.3950 1.3760 1.3184 1.4044
	6.4050 1.3560 1.2876 1.3972
	6.4150 1.3540 1.2956 1.4128
	6.4250 1.3480 1.2920 1.3960
	6.4350 1.3560 1.2668 1.4004
	6.4450 1.3540 1.2712 1.3960
	6.4550 1.3420 1.2996 1.3884
	6.4650 1.3380 1.2992 1.3840
	6.4750 1.3380 1.2876 1.3776
	6.4850 1.3320 1.2936 1.3896
	6.4950 1.3400 1.2956 1.3720
	6.5050 1.3520 1.3104 1.4040
	6.5150 1.3700 1.2880 1.4204
	6.5250 1.3440 1.2916 1.4084
	6.5350 1.3440 1.2796 1.3888
	6.5450 1.3340 1.2720 1.4012
	6.5550 1.3180 1.2536 1.3888
	6.5650 1.3320 1.2836 1.3732
	6.5750 1.3740 1.2996 1.4004
	6.5850 1.3360 1.2632 1.3880
	6.5950 1.3420 1.2916 1.3920
	6.6050 1.3560 1.2904 1.4088
	6.6150 1.3260 1.2664 1.3844
	6.6250 1.3300 1.2628 1.3764
	6.6350 1.3300 1.2920 1.3812
	6.6450 1.3480 1.2960 1.4040
	6.6550 1.3320 1.2920 1.4000
	6.6650 1.3600 1.2868 1.4164
	6.6750 1.3460 1.3064 1.3884
	6.6850 1.3500 1.2672 1.3968
	6.6950 1.3660 1.3200 1.4084
	6.7050 1.3200 1.2672 1.3960
	6.7150 1.3580 1.3076 1.4124
	6.7250 1.3420 1.3068 1.3768
	6.7350 1.3440 1.2720 1.4004
	6.7450 1.3500 1.2596 1.4084
	6.7550 1.3260 1.2712 1.3880
	6.7650 1.3240 1.2672 1.3884
	6.7750 1.3420 1.2864 1.3776
	6.7850 1.3460 1.2756 1.3848
	6.7950 1.3560 1.3040 1.4164
	6.8050 1.3440 1.2868 1.3884
	6.8150 1.3340 1.2676 1.4172
	6.8250 1.3480 1.2876 1.3920
	6.8350 1.3480 1.2800 1.3932
	6.8450 1.3440 1.3056 1.4004
	6.8550 1.3560 1.3152 1.3928
	6.8650 1.3440 1.2916 1.3848
	6.8750 1.3400 1.2676 1.3964
	6.8850 1.3520 1.2944 1.3884
	6.8950 1.3540 1.2920 1.3980
	6.9050 1.3260 1.2636 1.3720
	6.9150 1.3520 1.2908 1.3960
	6.9250 1.3660 1.3076 1.4000
	6.9350 1.3560 1.3076 1.4080
	6.9450 1.3400 1.2924 1.3940
	6.9550 1.3340 1.2712 1.3844
	6.9650 1.3460 1.3064 1.4040
	6.9750 1.3400 1.2680 1.3764
	6.9850 1.3400 1.2796 1.3928
	6.9950 1.3560 1.3068 1.4012
	7.0050 1.3400 1.2796 1.3880
	7.0150 1.3380 1.2680 1.3936
	7.0250 1.3300 1.2636 1.3600
	7.0350 1.3460 1.2840 1.3960
	7.0450 1.3580 1.3000 1.4128
	7.0550 1.3700 1.3040 1.4208
	7.0650 1.3760 1.3076 1.4208
	7.0750 1.3560 1.3136 1.4040
	7.0850 1.3680 1.2796 1.4280
	7.0950 1.3520 1.2840 1.3840
	7.1050 1.3420 1.2756 1.3880
	7.1150 1.3560 1.2992 1.4008
	7.1250 1.3460 1.2756 1.4052
	7.1350 1.3320 1.2836 1.3732
	7.1450 1.3480 1.2792 1.3880
	7.1550 1.3500 1.3000 1.4040
	7.1650 1.3340 1.2876 1.3772
	7.1750 1.3360 1.2716 1.3720
	7.1850 1.3480 1.3040 1.4056
	7.1950 1.3500 1.2948 1.3884
	7.2050 1.3660 1.3312 1.4040
	7.2150 1.3100 1.2752 1.3804
}\rwItchAll

\pgfplotstableread{
	x med lo hi
	0.2319 1.4000 1.3636 1.4524
	0.4239 1.4160 1.3672 1.4564
	0.5993 1.4060 1.3720 1.4600
	0.7925 1.4060 1.3792 1.4408
	1.0344 1.4160 1.3796 1.4328
	1.2922 1.4140 1.3852 1.4524
	1.5139 1.4260 1.3680 1.4520
	1.7118 1.3960 1.3596 1.4444
	1.8883 1.3840 1.3276 1.4232
	1.9730 1.3840 1.3560 1.4164
	2.1195 1.3920 1.3552 1.4520
	2.3277 1.4000 1.3388 1.4400
	2.4770 1.3940 1.3400 1.4288
	2.5931 1.4000 1.3636 1.4484
	2.7303 1.4040 1.3596 1.4440
	2.8811 1.4000 1.3640 1.4480
	3.0688 1.4160 1.3720 1.4524
	3.2074 1.4360 1.3880 1.4600
	3.2970 1.4120 1.3788 1.4440
	3.3486 1.3800 1.3220 1.4084
	3.3812 1.3740 1.2432 1.4244
	3.4093 1.4060 1.2916 1.4684
	3.4656 1.4140 1.3704 1.4400
	3.4890 1.4100 1.3588 1.4404
	3.5114 1.4280 1.3956 1.4524
	3.5350 1.4140 1.3796 1.4480
	3.5589 1.4100 1.3832 1.4400
	3.5836 1.4100 1.3796 1.4408
	3.6094 1.4160 1.3800 1.4484
	3.6358 1.4100 1.3676 1.4404
	3.6621 1.4180 1.3760 1.4484
	3.6850 1.4120 1.3880 1.4444
	3.7097 1.4000 1.3676 1.4408
	3.7294 1.4200 1.3760 1.4400
	3.7524 1.4240 1.3840 1.4484
	3.7750 1.4200 1.3920 1.4440
	3.7987 1.4320 1.3960 1.4564
	3.8207 1.4140 1.3796 1.4440
	3.8450 1.4220 1.3676 1.4440
	3.8709 1.4240 1.4000 1.4528
	3.8999 1.4160 1.3996 1.4484
	3.9296 1.4220 1.3800 1.4520
	3.9536 1.4160 1.3880 1.4480
	3.9788 1.4240 1.3828 1.4524
	4.0061 1.4140 1.3712 1.4400
	4.0278 1.4200 1.3712 1.4448
	4.0508 1.4360 1.4040 1.4640
	4.0748 1.4360 1.4000 1.4524
	4.0971 1.4160 1.3760 1.4560
	4.1253 1.4200 1.3876 1.4440
	4.1475 1.4140 1.3956 1.4364
	4.1691 1.4240 1.3960 1.4484
	4.1901 1.4240 1.3800 1.4488
	4.2083 1.4240 1.4032 1.4520
	4.2263 1.4320 1.4028 1.4520
	4.2443 1.4280 1.4036 1.4524
	4.2699 1.4440 1.4200 1.4640
	4.2931 1.4320 1.3996 1.4488
	4.3149 1.4360 1.4076 1.4560
	4.3338 1.4240 1.3880 1.4528
	4.3530 1.4200 1.3916 1.4452
	4.3744 1.4240 1.3960 1.4560
	4.3945 1.4140 1.3840 1.4480
	4.4157 1.4320 1.3952 1.4560
	4.4375 1.4160 1.3840 1.4488
	4.4583 1.4240 1.3920 1.4484
	4.4755 1.4280 1.3960 1.4600
	4.4954 1.4220 1.3876 1.4640
	4.5188 1.4060 1.3840 1.4444
	4.5412 1.4140 1.3832 1.4480
	4.5631 1.4200 1.3920 1.4564
	4.5822 1.4240 1.3988 1.4484
	4.6079 1.4100 1.3876 1.4404
	4.6310 1.4080 1.3716 1.4404
	4.6523 1.4040 1.3680 1.4280
	4.6719 1.4160 1.3836 1.4404
	4.6943 1.4160 1.3872 1.4524
	4.7159 1.4220 1.3872 1.4560
	4.7407 1.4040 1.3640 1.4368
	4.7652 1.4020 1.3676 1.4320
	4.7881 1.4000 1.3716 1.4440
	4.8132 1.4020 1.3596 1.4404
	4.8368 1.4120 1.3880 1.4404
	4.8565 1.4180 1.3744 1.4404
	4.8786 1.4100 1.3672 1.4560
	4.9022 1.3960 1.3628 1.4280
	4.9266 1.4060 1.3752 1.4324
	4.9451 1.4020 1.3508 1.4480
	4.9646 1.4120 1.3748 1.4484
	4.9885 1.4060 1.3668 1.4360
	5.0103 1.4200 1.3760 1.4484
	5.0341 1.4180 1.3920 1.4488
	5.0598 1.4140 1.3760 1.4560
	5.0867 1.4260 1.3876 1.4604
	5.1117 1.4100 1.3792 1.4368
	5.1396 1.4120 1.3556 1.4480
	5.1679 1.4160 1.3840 1.4524
	5.1961 1.4140 1.3716 1.4484
	5.2242 1.4060 1.3632 1.4444
	5.2529 1.4080 1.3800 1.4400
	5.2781 1.4120 1.3800 1.4400
	5.3076 1.4080 1.3836 1.4404
	5.3352 1.4180 1.3840 1.4444
	5.3592 1.4180 1.3840 1.4600
	5.3845 1.4140 1.3876 1.4444
	5.4116 1.4160 1.3796 1.4448
	5.4400 1.4100 1.3864 1.4444
	5.4641 1.4220 1.3796 1.4480
	5.4883 1.4060 1.3840 1.4440
	5.5177 1.4060 1.3708 1.4440
	5.5451 1.4160 1.3704 1.4440
	5.5669 1.4100 1.3716 1.4444
	5.5940 1.4120 1.3796 1.4408
	5.6241 1.4040 1.3556 1.4440
	5.6463 1.4020 1.3396 1.4520
	5.6735 1.4120 1.3508 1.4444
	5.6965 1.4160 1.3820 1.4404
	5.7238 1.4080 1.3640 1.4480
	5.7536 1.4080 1.3756 1.4372
	5.7752 1.4100 1.3720 1.4408
	5.7996 1.4160 1.3796 1.4480
	5.8268 1.4060 1.3436 1.4484
	5.8496 1.4120 1.3840 1.4440
	5.8734 1.4160 1.3600 1.4560
	5.9010 1.4160 1.3760 1.4480
	5.9288 1.4040 1.3596 1.4440
	5.9580 1.4160 1.3832 1.4524
	5.9833 1.4140 1.3676 1.4520
	6.0180 1.4160 1.3920 1.4480
	6.0424 1.4320 1.4000 1.4564
	6.0704 1.4200 1.3680 1.4524
	6.0980 1.4100 1.3756 1.4440
	6.1280 1.4000 1.3600 1.4400
	6.1508 1.4080 1.3800 1.4560
	6.1779 1.4140 1.3716 1.4484
	6.2082 1.4120 1.3680 1.4440
	6.2353 1.4160 1.3756 1.4564
	6.2623 1.4040 1.3876 1.4360
	6.2878 1.4040 1.3556 1.4360
	6.3166 1.4140 1.3720 1.4404
	6.3438 1.4240 1.3880 1.4564
	6.3704 1.4180 1.3756 1.4520
	6.3945 1.4220 1.3880 1.4520
	6.4202 1.4240 1.3876 1.4560
	6.4463 1.4200 1.3868 1.4480
	6.4708 1.4080 1.3788 1.4284
	6.4965 1.4020 1.3600 1.4364
	6.5200 1.4160 1.3876 1.4484
	6.5452 1.4080 1.3704 1.4404
	6.5715 1.4080 1.3708 1.4440
	6.5979 1.4080 1.3800 1.4360
	6.6259 1.3960 1.3628 1.4320
	6.6475 1.4000 1.3720 1.4328
	6.6712 1.4080 1.3836 1.4364
	6.6937 1.4240 1.3832 1.4564
	6.7204 1.4020 1.3760 1.4284
	6.7483 1.3960 1.3680 1.4408
	6.7804 1.4120 1.3668 1.4360
	6.8038 1.4120 1.3800 1.4440
	6.8249 1.4040 1.3756 1.4416
	6.8497 1.4120 1.3516 1.4520
	6.8751 1.4140 1.3720 1.4440
	6.9025 1.4160 1.3800 1.4524
	6.9301 1.4160 1.3828 1.4444
	6.9631 1.4120 1.3760 1.4440
	6.9931 1.4080 1.3676 1.4440
	7.0233 1.4000 1.3712 1.4248
	7.0500 1.4180 1.3832 1.4520
	7.0735 1.4200 1.3840 1.4524
	7.0996 1.4200 1.3840 1.4440
	7.1236 1.3980 1.3476 1.4440
	7.1498 1.3860 1.3504 1.4084
	7.1803 1.3840 1.3600 1.4240
	7.2051 1.3960 1.3748 1.4284
}\rwItchQQQ

\begin{figure}[!t]
	\centering
	{%
		\begin{tikzpicture}[trim axis right]
			\begin{axis}[width=\columnwidth, height=0.5\columnwidth, axis lines=box, axis line style={black, line width=0.5pt}, tick align=inside, tick style={black!70},
					grid=both, minor x tick num=1, minor y tick num=1,
					major grid style={black!40, dotted, line width=0.5pt}, minor grid style={black!18, dotted, line width=0.4pt},
					tick label style={font=\scriptsize}, label style={font=\footnotesize}, title style={font=\footnotesize},
					ylabel={Cycles per Operation}, ylabel style={yshift=-2pt},
					legend style={font=\scriptsize, draw=black!60, fill=white, fill opacity=0.9, text opacity=1, legend cell align=left},
					xmin=0, xmax=7.25, ymin=0.95, ymax=1.65, ytick={1.0,1.2,1.4,1.6},
					xlabel={Operations (Millions)}, title={NASDAQ ITCH order flow},
					legend style={at={(0.98,0.97)}, anchor=north east, legend columns=2, column sep=5pt}]
				\draw[black!60, line width=0.6pt, loosely dashed] (axis cs:0,1) -- (axis cs:7.25,1);
				\addplot[draw=none, name path=ialo, forget plot] table[x=x, y=lo]{\rwItchAll};
				\addplot[draw=none, name path=iahi, forget plot] table[x=x, y=hi]{\rwItchAll};
				\addplot[fill=black!18, fill opacity=0.35, draw=none, forget plot] fill between[of=ialo and iahi];
				\addplot[black, line width=1.0pt, forget plot] table[x=x, y=med]{\rwItchAll};
				\addplot[draw=none, name path=iqlo, forget plot] table[x=x, y=lo]{\rwItchQQQ};
				\addplot[draw=none, name path=iqhi, forget plot] table[x=x, y=hi]{\rwItchQQQ};
				\addplot[fill=black!10, fill opacity=0.35, draw=none, forget plot] fill between[of=iqlo and iqhi];
				\addplot[black!60, line width=1.0pt, densely dashed, forget plot] table[x=x, y=med]{\rwItchQQQ};
				\addlegendimage{black, line width=1.0pt}\addlegendentry{All eight symbols}
				\addlegendimage{black!60, line width=1.0pt, densely dashed}\addlegendentry{QQQ alone}
			\end{axis}
		\end{tikzpicture}%
	}
	\caption{Cycles per operation of a pipeline that stalls whenever an operation needs a write still in flight, over NASDAQ ITCH order flow for all eight symbols and for QQQ alone. The dashed line marks one cycle per operation; lines are medians, and bands span the 10th to 90th percentiles.}
	\label{fig:rw_throughput}
\end{figure}

\subsection{Pipelined Reads, Writes and Their Hazards}\label{sec:mot_rw}

A tough problem to engineer systems around is one in which several reads and
writes are in flight at once, all destined for the same nodes. An event in an
order book, for example, changes a price level and is followed by a query about
the book's depth, so reads and writes arrive in equal numbers.

Fig.~\ref{fig:rw_throughput} quantifies what this costs a pipeline that could
issue one operation every cycle. Here, every subsequent operation in need of an
unfinished write must stall until the write is complete. The trace is NASDAQ
ITCH order flow for eight symbols: 3.6 million order-book updates, each
followed by a depth query at the best price of the side just updated. On it,
the stalling pipeline needs 1.33 cycles per operation across the eight
interleaved symbols, and 1.41 when it serves the book of the busiest symbol,
QQQ, alone, because a single book crowds its updates onto one price ladder.

A system that is to keep pace with such a workload must therefore accept a
read and a write in the same stream every cycle, and \emph{handle hazards
	gracefully}: neither padded with no-op bubbles nor fraught with stalls, which
here cost 33\% to 41\% more cycles per operation.

\begin{figure*}[!t]
	\centering
	{%
		\definecolor{rS}{HTML}{2B6CB0}\definecolor{rX}{HTML}{C0392B}\definecolor{cLn}{HTML}{E4572E}%
		\pgfplotsset{cS1/.style={rS!50, line width=0.8pt}, cS2/.style={rS, line width=0.8pt, densely dashed},
			cI1/.style={black, line width=0.9pt}, cI2/.style={black!55, line width=0.9pt, dash dot},
			cX1/.style={rX!60, line width=0.8pt, dotted}, cX2/.style={rX, line width=0.9pt}}%
		\begin{tikzpicture}
			\begin{axis}[width=0.42\textwidth, height=0.25\textwidth, axis lines=box, axis line style={black, line width=0.5pt}, tick align=inside, tick style={black!70},
					xmode=log, log basis x=2, ymode=log, grid=both,
					major grid style={black!40, dotted, line width=0.5pt}, minor grid style={black!15, dotted, line width=0.4pt},
					tick label style={font=\scriptsize}, label style={font=\footnotesize}, title style={font=\footnotesize},
					xlabel={Latency $L$ (Cycles)}, ylabel={Share Slower than $L$}, ylabel style={yshift=-4pt},
					ymin=1e-05, ymax=1.6, xmin=3.6, xmax=12000,
					xtick={4,16,64,256,1024,4096}, xticklabels={$2^{2}$,$2^{4}$,$2^{6}$,$2^{8}$,$2^{10}$,$2^{12}$}, title={(a) Fenwick}]
				\addplot[cS1] coordinates {(4,0.4054) (5,0.01985) (6,0.001533) (7,3.3e-05) (8,1e-05)};
				\addplot[cS2] coordinates {(4,0.81575) (5,0.5119) (6,0.326917) (7,0.21175) (8,0.138983) (9,0.091683) (10,0.06295) (11,0.042567) (12,0.029117) (13,0.018617) (14,0.0105) (15,0.007) (16,0.0038) (17,0.001983) (18,0.0007) (19,0.00025) (20,1e-05)};
				\addplot[cI1] coordinates {(4,0.976217) (5,0.923233) (6,0.871867) (7,0.8199) (8,0.767783) (9,0.7197) (10,0.67475) (11,0.630467) (12,0.588383) (13,0.546383) (14,0.5018) (15,0.45995) (16,0.423133) (17,0.3917) (18,0.365517) (19,0.339467) (20,0.313683) (21,0.290233) (22,0.264767) (23,0.240167) (24,0.220283) (25,0.202283) (26,0.180067) (27,0.158767) (28,0.13945) (29,0.122317) (30,0.108467) (31,0.095333) (32,0.083567) (33,0.072783) (34,0.06075) (35,0.053117) (36,0.047617) (37,0.040917) (38,0.035983) (39,0.030633) (40,0.024883) (41,0.020167) (42,0.016267) (43,0.01255) (44,0.010183) (45,0.00775) (46,0.006067) (47,0.00475) (48,0.003667) (49,0.002833) (50,0.002417) (51,0.001833) (52,0.0014) (53,0.000833) (54,0.0003) (55,5e-05) (56,1e-05)};
				\addplot[cI2] coordinates {(4,0.995983) (5,0.989017) (6,0.983183) (7,0.976517) (8,0.969933) (9,0.964783) (10,0.959833) (11,0.95545) (12,0.9512) (13,0.946367) (14,0.940133) (15,0.933033) (16,0.925767) (17,0.919533) (18,0.912483) (19,0.906717) (20,0.901683) (21,0.8974) (22,0.892717) (23,0.8879) (24,0.883583) (25,0.87825) (26,0.871667) (27,0.86505) (28,0.85925) (29,0.853) (30,0.847133) (31,0.84265) (32,0.839683) (33,0.836617) (34,0.832783) (35,0.82925) (36,0.82485) (37,0.820917) (38,0.8166) (39,0.812317) (40,0.808517) (41,0.803383) (42,0.798083) (43,0.7941) (44,0.78935) (45,0.78595) (46,0.7803) (47,0.77625) (48,0.77115) (49,0.766067) (50,0.7632) (51,0.76085) (52,0.758883) (53,0.7568) (54,0.7541) (55,0.7518) (56,0.748267) (57,0.743) (58,0.738917) (59,0.7339) (60,0.728733) (61,0.723817) (62,0.719717) (63,0.716817) (64,0.713067) (65,0.70815) (66,0.70165) (67,0.697067) (68,0.692417) (69,0.68895) (71,0.68005) (72,0.674567) (73,0.669883) (74,0.666083) (75,0.66205) (77,0.653533) (78,0.6497) (79,0.644917) (81,0.633967) (82,0.628583) (83,0.624733) (85,0.61625) (86,0.61065) (87,0.604317) (89,0.58845) (90,0.580733) (92,0.564267) (93,0.55725) (95,0.544117) (97,0.535083) (98,0.529567) (100,0.523333) (101,0.520483) (103,0.516033) (105,0.511333) (107,0.508083) (108,0.506817) (110,0.5046) (112,0.50355) (114,0.50075) (116,0.498367) (118,0.495833) (120,0.49345) (122,0.491383) (124,0.490233) (126,0.489617) (128,0.4894) (130,0.489083) (132,0.488867) (134,0.488467) (137,0.487917) (139,0.485783) (141,0.482733) (143,0.47855) (146,0.472183) (148,0.466983) (151,0.45985) (153,0.454083) (156,0.4444) (158,0.4368) (161,0.422067) (164,0.408367) (166,0.397467) (169,0.3857) (172,0.370017) (175,0.347583) (178,0.321267) (181,0.29455) (184,0.264683) (187,0.241567) (190,0.227317) (193,0.216783) (196,0.20975) (200,0.189183) (203,0.174417) (206,0.164117) (210,0.15245) (213,0.141433) (217,0.132133) (220,0.122267) (224,0.108117) (228,0.096917) (231,0.090983) (235,0.08085) (239,0.07505) (243,0.067567) (247,0.058167) (251,0.0482) (256,0.0316) (260,0.016283) (264,0.010033) (265,0.008417) (266,0.0069) (267,0.005317) (268,0.003417) (269,0.00255) (270,0.0018) (271,0.000633) (272,0.000267) (273,1e-05)};
				\addplot[cX1] coordinates {(4,0.999467) (5,0.99885) (6,0.998217) (7,0.996867) (8,0.9961) (9,0.995233) (10,0.994833) (11,0.994367) (12,0.993867) (13,0.993533) (14,0.993167) (15,0.99285) (16,0.992233) (17,0.991133) (18,0.98995) (19,0.988683) (20,0.986433) (21,0.98435) (22,0.983567) (23,0.982083) (24,0.98075) (25,0.979667) (26,0.978883) (27,0.97835) (28,0.977967) (29,0.977767) (30,0.977333) (31,0.976683) (32,0.976167) (33,0.97585) (34,0.975617) (35,0.97525) (36,0.9752) (37,0.975183) (38,0.975183) (39,0.97515) (40,0.975083) (41,0.975083) (42,0.975083) (43,0.975083) (45,0.974983) (46,0.9749) (47,0.974517) (48,0.974133) (49,0.974067) (51,0.973433) (52,0.9729) (53,0.97235) (55,0.971733) (56,0.9712) (58,0.97035) (59,0.9699) (61,0.969217) (62,0.968817) (64,0.968483) (65,0.967933) (67,0.966817) (69,0.9666) (71,0.966183) (72,0.96585) (74,0.965433) (76,0.965183) (78,0.964917) (80,0.964583) (82,0.962933) (85,0.961517) (87,0.9612) (89,0.9611) (91,0.96105) (94,0.9599) (96,0.959467) (99,0.95785) (101,0.95705) (104,0.95555) (106,0.95445) (109,0.953917) (112,0.953467) (115,0.9523) (118,0.951867) (121,0.95) (124,0.94875) (127,0.9482) (131,0.947467) (134,0.947) (138,0.945783) (141,0.944883) (145,0.9437) (149,0.94315) (152,0.942683) (156,0.941417) (160,0.940017) (165,0.939767) (169,0.9386) (173,0.938033) (178,0.9361) (182,0.935733) (187,0.93435) (192,0.9328) (197,0.929433) (202,0.926667) (207,0.91835) (213,0.91455) (218,0.912033) (224,0.911083) (230,0.909383) (236,0.90835) (242,0.904483) (248,0.902333) (255,0.900567) (261,0.900033) (268,0.897317) (275,0.893333) (282,0.8907) (289,0.8887) (297,0.887017) (305,0.881017) (312,0.879517) (321,0.87675) (329,0.8741) (337,0.871167) (346,0.865883) (355,0.864517) (364,0.8613) (374,0.855783) (384,0.85195) (394,0.84995) (404,0.845917) (414,0.844167) (425,0.838883) (436,0.830483) (447,0.81515) (459,0.808467) (471,0.802083) (483,0.799083) (496,0.79645) (509,0.7934) (522,0.789383) (535,0.780683) (549,0.772783) (563,0.769483) (578,0.766083) (593,0.7627) (608,0.749117) (624,0.741667) (640,0.738433) (657,0.731717) (674,0.725833) (692,0.72185) (710,0.714633) (728,0.71225) (747,0.70855) (766,0.705583) (786,0.702) (807,0.6968) (828,0.685733) (849,0.682567) (871,0.6778) (894,0.67085) (917,0.66665) (941,0.660017) (965,0.650817) (990,0.6454) (1016,0.634467) (1042,0.621633) (1069,0.613717) (1097,0.60795) (1126,0.5936) (1155,0.590683) (1185,0.581433) (1216,0.565633) (1247,0.55405) (1280,0.546583) (1313,0.5379) (1347,0.52485) (1382,0.513417) (1418,0.50225) (1455,0.495667) (1492,0.484833) (1531,0.47215) (1571,0.464517) (1612,0.45205) (1654,0.438383) (1697,0.4134) (1741,0.39745) (1786,0.384833) (1832,0.365783) (1880,0.359267) (1929,0.339733) (1979,0.325867) (2030,0.306933) (2083,0.279933) (2137,0.26485) (2192,0.2433) (2249,0.232583) (2308,0.213417) (2368,0.17675) (2429,0.153133) (2492,0.135483) (2557,0.1189) (2623,0.100983) (2691,0.06985) (2761,0.037633) (2817,0.010217) (2818,0.00955) (2819,0.009217) (2820,0.008717) (2821,0.007983) (2822,0.00705) (2823,0.006183) (2824,0.00565) (2825,0.004933) (2826,0.003967) (2827,0.0037) (2828,0.003183) (2829,0.002233) (2830,0.001167) (2831,0.000267) (2832,8.3e-05) (2833,1e-05)};
				\addplot[cX2] coordinates {(4,0.999733) (5,0.999667) (6,0.999567) (7,0.999333) (8,0.999333) (9,0.999333) (10,0.999283) (11,0.999233) (12,0.99915) (13,0.998883) (14,0.998867) (15,0.998833) (16,0.998817) (17,0.998667) (18,0.998617) (19,0.998567) (20,0.998567) (21,0.998567) (22,0.998517) (23,0.998517) (24,0.9984) (25,0.998333) (26,0.998317) (27,0.997933) (28,0.997767) (29,0.997567) (30,0.997533) (31,0.997517) (32,0.997367) (33,0.99735) (34,0.997233) (35,0.997167) (36,0.997133) (37,0.9971) (38,0.997) (39,0.996983) (41,0.99675) (42,0.9966) (43,0.996367) (44,0.996233) (46,0.99615) (47,0.99585) (49,0.9956) (50,0.9956) (52,0.995483) (53,0.9954) (55,0.995383) (57,0.995317) (58,0.995133) (60,0.995) (62,0.994833) (64,0.994683) (66,0.99435) (68,0.9941) (70,0.993933) (72,0.993783) (75,0.9935) (77,0.9934) (79,0.993333) (82,0.993083) (84,0.99295) (87,0.9929) (90,0.992833) (92,0.99275) (95,0.992583) (98,0.992083) (101,0.99205) (104,0.991967) (108,0.9916) (111,0.991317) (114,0.9909) (118,0.990167) (121,0.9898) (125,0.9892) (129,0.988917) (133,0.98825) (137,0.987933) (141,0.987683) (146,0.987267) (150,0.9865) (155,0.985917) (160,0.985633) (165,0.984933) (170,0.984433) (175,0.984033) (181,0.982933) (186,0.982717) (192,0.982267) (198,0.981567) (204,0.980717) (210,0.979917) (217,0.979233) (223,0.97895) (230,0.978283) (238,0.9773) (245,0.976983) (252,0.976633) (260,0.9756) (268,0.974883) (277,0.974067) (285,0.973567) (294,0.972617) (303,0.971683) (313,0.9709) (322,0.970483) (332,0.969667) (342,0.968617) (353,0.9679) (364,0.9668) (375,0.965667) (387,0.964317) (399,0.963233) (411,0.96215) (424,0.9609) (437,0.9599) (450,0.959183) (464,0.957717) (479,0.9559) (494,0.954667) (509,0.953133) (525,0.952) (541,0.950383) (558,0.948133) (575,0.94685) (593,0.945083) (611,0.94355) (630,0.941233) (649,0.939267) (670,0.937683) (690,0.93625) (712,0.932983) (734,0.931533) (756,0.929533) (780,0.92705) (804,0.92585) (829,0.923817) (854,0.9216) (881,0.919333) (908,0.91685) (936,0.9146) (965,0.912367) (995,0.909883) (1026,0.90735) (1058,0.904317) (1090,0.902117) (1124,0.89985) (1159,0.896333) (1195,0.8929) (1232,0.8898) (1270,0.886483) (1309,0.882083) (1350,0.87845) (1391,0.875533) (1434,0.87095) (1479,0.867267) (1525,0.862517) (1572,0.858167) (1620,0.854) (1671,0.849167) (1722,0.843317) (1776,0.83795) (1830,0.83185) (1887,0.827333) (1946,0.822133) (2006,0.815767) (2068,0.810617) (2132,0.806217) (2198,0.79765) (2266,0.792217) (2336,0.7854) (2408,0.77745) (2483,0.770583) (2560,0.763933) (2639,0.755017) (2720,0.748383) (2805,0.739283) (2891,0.729483) (2981,0.720967) (3073,0.711883) (3168,0.70145) (3266,0.692117) (3367,0.68295) (3472,0.6724) (3579,0.662517) (3690,0.650267) (3804,0.640933) (3922,0.6316) (4043,0.61895) (4168,0.60585) (4297,0.594) (4430,0.579933) (4567,0.56365) (4709,0.551733) (4854,0.538633) (5005,0.52475) (5159,0.509283) (5319,0.493067) (5484,0.47945) (5653,0.46415) (5828,0.44695) (6009,0.430817) (6195,0.4123) (6386,0.39505) (6584,0.375883) (6788,0.35565) (6998,0.336767) (7214,0.315683) (7438,0.29655) (7668,0.276467) (7905,0.255467) (8150,0.232683) (8402,0.211817) (8662,0.1866) (8930,0.159217) (9206,0.120167) (9491,0.075183) (9768,0.0102) (9769,0.009467) (9770,0.00945) (9771,0.00925) (9772,0.009167) (9773,0.007783) (9774,0.007683) (9775,0.007617) (9776,0.0076) (9777,0.00725) (9778,0.0061) (9779,0.0059) (9780,0.0054) (9781,0.005183) (9782,0.003) (9783,0.0021) (9784,0.000917) (9785,1e-05)};
				\draw[black!55, line width=0.5pt, densely dashed] (axis cs:8,1e-5) -- (axis cs:8,1.6);
				\draw[black!55, line width=0.5pt, densely dashed] (axis cs:3.6,1e-3) -- (axis cs:12000,1e-3);
				\node[font=\tiny, text=black!70, anchor=south west, inner sep=1pt, rotate=90] at (axis cs:8.4,1.5e-5) {8 cycles};
				\node[font=\tiny, text=black!70, anchor=south west, inner sep=1pt] at (axis cs:3.8,1.12e-3) {1 in 1{,}000};
			\end{axis}
		\end{tikzpicture}\hspace{0.06\textwidth}
		\begin{tikzpicture}
			\begin{axis}[width=0.42\textwidth, height=0.25\textwidth, axis lines=box, axis line style={black, line width=0.5pt},
					tick align=inside, tick style={black!70}, xmin=0, xmax=1, ymin=3.4, ymax=40000, clip=false,
					ymode=log, log basis y=2, xtick={0,0.2909,0.5818,0.7714,1}, xticklabels={0,0.4,0.8,0.9,1},
					ytick={4,64,1024,16384}, yticklabels={$2^{2}$,$2^{6}$,$2^{10}$,$2^{14}$},
					grid=both, major grid style={black!40, dotted, line width=0.5pt}, minor grid style={black!15, dotted, line width=0.4pt},
					tick label style={font=\scriptsize}, label style={font=\footnotesize}, title style={font=\footnotesize, yshift=6pt},
					xlabel={Offered Load}, ylabel={Median (Cycles)}, ylabel style={yshift=-4pt}, title={(b) Median latency by load}]
				\draw[black!55, line width=0.6pt, dash dot] (axis cs:0,4) -- (axis cs:1,4);
				\node[font=\tiny, text=black!70, anchor=south west, inner sep=1.5pt] at (axis cs:0.0145,4.15) {Conflict-free};
				\addplot[black!45, line width=0.7pt] coordinates {(0.0364,4.0) (0.0727,4.0) (0.1091,4.0) (0.1455,4.0) (0.1818,4.0) (0.2182,4.0) (0.2545,4.0) (0.2909,4.0) (0.3273,4.0) (0.3636,5.0) (0.4000,5.0) (0.4364,5.0) (0.4727,5.0) (0.5091,5.0) (0.5455,5.0) (0.5818,6.0) (0.6114,7.0) (0.6343,8.0) (0.6571,9.0) (0.6800,15.0) (0.7029,115.0) (0.7257,518.0) (0.7486,749.0) (0.7714,1422.0) (0.7943,1748.0) (0.8171,2068.0) (0.8400,2621.0) (0.8629,2936.0) (0.8857,3352.0) (0.9086,3668.0) (0.9314,4232.0) (0.9543,4436.5) (0.9771,5075.5) (1.0000,5250.0)};
				\draw[cLn, densely dashed, line width=0.8pt] (axis cs:0.6914285714285715,3.4) -- (axis cs:0.6914285714285715,40000);
				\node[font=\scriptsize, text=cLn, anchor=south, inner sep=1.5pt] at (axis cs:0.6914285714285715,40000) {0.865};
				\addplot[only marks, mark=*, mark size=1.6pt, draw=white, line width=0.3pt, fill=rS!50] coordinates {(0.1455,4.0)}; \addplot[only marks, mark=*, mark size=1.6pt, draw=white, line width=0.3pt, fill=rS] coordinates {(0.5818,6.0)}; \addplot[only marks, mark=*, mark size=1.6pt, draw=white, line width=0.3pt, fill=black] coordinates {(0.6800,15.0)}; \addplot[only marks, mark=*, mark size=1.6pt, draw=white, line width=0.3pt, fill=black!55] coordinates {(0.7029,115.0)}; \addplot[only marks, mark=*, mark size=1.6pt, draw=white, line width=0.3pt, fill=rX!60] coordinates {(0.7714,1422.0)}; \addplot[only marks, mark=*, mark size=1.6pt, draw=white, line width=0.3pt, fill=rX] coordinates {(1.0000,5250.0)};
			\end{axis}
		\end{tikzpicture}%
	}
	\caption{Fenwick query latency under XOR-folded placement through the PULP TCDM crossbar: (a)~share of queries slower than $L$ at six loads; (b)~median latency against load, with the inflection and the loads drawn in (a) marked. The dash-dotted line marks a conflict-free placement, which completes every query in 4 cycles at any load.}
	\label{fig:conflict_ccdf}
\end{figure*}

\subsection{Desirable Properties of Placement Functions}\label{sec:mot_conflict}

What makes a placement good comes down to two properties: \emph{conflict
	freedom} and \emph{balance}.

\subsubsection{Conflict Freedom}

An operation on a tree reads a fixed set of nodes together, its
\emph{template}~\cite{das2000mappings}, whose size is the $d$ of
Section~\ref{sec:mot_banks}; here, a leaf-to-root path, as it arises in
Fenwick trees, segment trees and succinct tries. Two nodes of one template
that share a bank are in \emph{conflict}, and one must wait for a later cycle;
a placement is conflict-free when every template lies in distinct banks, so
that all its nodes can be read in the same cycle.

To see how much this matters in practice, we ran queries on a dense Fenwick
tree with 16-bit keys through the PULP TCDM
interconnect~\cite{pulp_tcdm,pulp_cluster_interconnect}. It is a full crossbar
that resolves conflicts at run time, with a round-robin arbiter in front of
each of sixteen one-cycle banks. We placed the tree by XOR folding, a linear
form of address hashing that balances the banks but does not keep a template
apart (cf.~Theorem~\ref{thm:linear}). Each query presents all of its nodes at
once, one per request port, and its latency runs from its arrival to its last
response, counted through the front end and the adder tree, so that a query
that meets no conflict takes four cycles. We swept the offered load from 0.01
to 1.00 queries per cycle and drew, at six loads, the share of queries slower
than each latency (Fig.~\ref{fig:conflict_ccdf}(a)), and the median latency at
every load (Fig.~\ref{fig:conflict_ccdf}(b)).

The latency divides into two regions at a knee
(Fig.~\ref{fig:conflict_ccdf}(b)), at a load of 0.865 queries per cycle. Below
the knee, the bulk of the queries finishes within a few cycles, with a short
and thin tail; at a load of 0.2, the slowest query in a thousand takes
7 cycles. Approaching the knee, the tail lengthens and fattens together: the
99.9th percentile rises sixfold between loads of 0.80 and 0.86, from 19 cycles
to 122.

Beyond the knee, the distribution no longer has a tail in the usual sense: it
has moved as a whole. The queue never drains, so the median itself climbs, to
	5{,}273 cycles at full load, and fewer than one query in a thousand finishes in
the pipeline's four cycles. A placement that balances the banks but lets
templates collide thus caps throughput well short of one query per cycle.

\subsubsection{Balance}
Balance has two senses. \emph{Capacity} balance is how many nodes each bank
holds: it belongs to the placement alone, is fixed when the tree is laid out,
and decides how large every bank must be built, since all are sized for the
fullest. \emph{Access} balance is how often each bank is touched: it depends on
the workload as well, and decides how hard each bank works.

Jiang et al.~\cite{jiang2026scamp} make the cost of capacity imbalance precise
for IPv6 lookup pipelines, the design now standard in routing ASICs and
RMT-style programmable switches: for a fixed memory budget $C$ and a per-entry
cost $\eta_{\mathrm{pp}}$, the number of entries a pipeline can hold is
$$
	N_{\mathrm{max}} = \frac{\eta_{\mathrm{util}} \cdot C}{\eta_{\mathrm{pp}}}.
$$
When every bank is provisioned for the fullest one, the utilisation is set by
the placement alone,
$$
	\eta_{\mathrm{util}} = \frac{\sum_b n_b}{B \cdot \max_b n_b},
$$
where $n_b$ is the number of nodes in bank $b$, so $N_{\mathrm{max}}$ falls in
direct proportion to the share of the fullest bank. Placing the 16-bit Fenwick
tree by level fills one bank with half of its nodes, so
$\eta_{\mathrm{util}} = 0.125$, and the same memory holds one eighth of the
entries an exactly balanced placement would; trace-derived
colouring~\cite{zhou2017new} sizes every bank for its heaviest load of 50,800
nodes, for $\eta_{\mathrm{util}} = 0.081$. Sparse structures inherit the
problem from their keys: in a forest of dense subtrees whose live subtrees
follow the pages of a real allocator trace (\texttt{git}), banking each subtree
by the low bits of its address reaches only $\eta_{\mathrm{util}} = 0.69$.

Access imbalance costs differently. Behind a cache, a bank asked to serve more
of the nodes a workload reads than its share evicts more of them, so access
imbalance turns into evictions; Section~\ref{sec:eval_c2} measures this on the
succinct tries of C$^2$~\cite{zhang2027c2}.
